\documentclass[10pt,a4paper]{amsart}
\usepackage[margin=19mm]{geometry}
\usepackage{amsmath,amssymb,mathtools}
\usepackage[scr=rsfs,cal=euler]{mathalfa}
\usepackage{xfrac}
\usepackage[unicode,hidelinks,bookmarks=false]{hyperref}
\newcommand{\ie}{i.e.\ }
\newcommand{\cf}{cf.\ }
\newcommand{\resp}{respectively\ }
\newcommand{\Subsection}{\subsection}
\providecommand{\nobreakdash}{\nobreak\hbox{-}}
\setlength{\emergencystretch}{2em}
\multlinegap0pt
\usepackage{amssymb}
\newtheorem{proposition}{Proposition}
\def\d{\mathsf d}
\def\r{\mathbb R}
\title{Generalization of the catenary in the Dual plane}
\author{Muhittin Evren Aydın, Rafael L\'opez}
\date{Source: arXiv:2603.02107v1 (CC BY 4.0)}
\begin{document}
\maketitle

\noindent\emph{Source and licence.} Generalization of the catenary in the Dual plane. Version arXiv:2603.02107v1 (CC BY 4.0), distributed by arXiv under the Creative Commons Attribution 4.0 International licence, http://creativecommons.org/licenses/by/4.0/, as recorded in the arXiv licence metadata for this exact version and shown on the abstract page https://arxiv.org/abs/2603.02107v1. Full licence notice: \texttt{licenses/arxiv-2603.02107-CC-BY-4.0.txt}.\par\smallskip

\noindent\textit{Editorial scope.} Counted unit: the article's proposition p25 (source0.tex lines 316--326). The article's complete original proof of it is delivered with it (source0.tex lines 328--361, one \texttt{proof} environment of 34 lines). No mathematics is written, repaired or re-derived: the statement and the proof are the article's own lines, in the article's own notation, and no retained line is substituted or re-flowed.  Retained uncounted as the source-local context the proof needs: lines 187--193; lines 201--203; lines 223--225.  Nothing was omitted from any retained span, so the package carries no omission record.  No retained line contains a citation, so the wrapper carries no bibliography and no bibliography entry.  No retained line contains a figure environment, an image-inclusion command or drawing code, so no image is delivered and the package declares no asset dependency.  Editorial formatting only, in addition: an AMS \texttt{amsart} preamble that restates the article's own declarations and macros used by the retained lines, namely the environments \texttt{proposition} on separate counters, together with the standard packages those lines need.  The retained environments print in this document as Proposition 1 in order of appearance, whereas the article prints the same items with the numbering of its own full document.  The complete native source of the article as distributed in the arXiv e-print for this exact version is bundled unchanged as \texttt{sources/195599/source0.tex}.
\begin{equation}\label{g}
\begin{split}
\gamma(x)&=\psi(x)+\varepsilon \phi(x),\\
\psi(x)&=(x,y(x)),\\
\phi(x)&=(w(x),z(x)),
\end{split}
\end{equation}

\begin{equation}\label{orto}
w'+y'z'=0.
\end{equation}

\begin{equation}\label{eq1}
 z'' + \alpha \frac{y'}{y}(z'+v) + \alpha \frac{z+vx}{y^2} = 0.
\end{equation}

\begin{proposition}\label{p25}
A curve $\gamma $ in $\d^2$     given by \eqref{g} is a $(-1)$-catenary if and only if 
\begin{equation}
\begin{split}\label{s2}
y(x)&=\sqrt{R^2-(x-m)^2},\\
w(x)&=(v-d_1) y(x) + d_2 (x-m) - d_2 y(x) \arcsin\left(\frac{x-m}{R}\right) + d_3,\\
z(x)&=-vx + d_1(x-m) + d_2 \left( y(x) + (x-m)\arcsin\left(\frac{x-m}{R}\right) \right),
\end{split}
\end{equation}
where $R, m, v, d_i \in \mathbb{R}$ with $R>0$. Furthermore, the maximal domain is given by $|x-m| < R$.
\end{proposition}

\begin{proof}
For $\alpha=-1$, the   curve $\psi$ is a $(-1)$-catenary in $\r^2$, which corresponds to a circular arc. Thus, we have $y(x) = \sqrt{R^2-(x-m)^2}$, and its derivative is $y'(x) = \frac{-(x-m)}{y(x)}$, where $m,R\in \r$, $R>0$. Substituting $\alpha=-1$ into \eqref{eq1}, we have
$$
z'' - \frac{y'}{y}(z'+v) - \frac{z+vx}{y^2} = 0.
$$
We introduce the change of variable $u(x)=z(x)+vx$, obtaining
$$
u'' - \frac{y'}{y} u' - \frac{1}{y^2} u = 0.
$$
Substituting $y$ and $y'$ into this equation yields
$$
u'' + \frac{x-m}{R^2-(x-m)^2} u' - \frac{1}{R^2-(x-m)^2} u = 0,
$$
or equivalently,
$$
(R^2-(x-m)^2) u'' + (x-m) u' - u = 0.
$$
It is straightforward to verify that $u_1(x) = x-m$ is a particular solution to this equation. A second, linearly independent solution can be found by using reduction of order, yielding 
$$u_2(x) = \sqrt{R^2-(x-m)^2} + (x-m)\arcsin\left(\frac{x-m}{R}\right).$$
 Noting that the first term of $u_2$ is exactly $y(x)$, the general solution for $u(x)$ is
$$
u(x) = d_1 (x-m) + d_2 \left( y(x) + (x-m)\arcsin\left(\frac{x-m}{R}\right) \right),
$$
where $d_1, d_2 \in \mathbb{R}$. Reverting the change of variable, we obtain the expression for $z(x)$ in  \eqref{s2}. 


Finally, we need to find $w$ from \eqref{orto}. Differentiating $z(x)$ gives
$$
z'(x) = -v + d_1 + d_2 \arcsin\left(\frac{x-m}{R}\right).
$$
Hence
$$w'(x) = \frac{x-m}{\sqrt{R^2-(x-m)^2}} \left( -v + d_1 + d_2 \arcsin\left(\frac{x-m}{R}\right) \right).$$
An integration yields the expression for $w(x)$ in \eqref{s2}.   This completes the proof.
\end{proof}

\end{document}
